\chapter{Annexes}

Le code source complet du projet est disponible à l'adresse~: \texttt{https://github.com/anomalyco/supervision\_thermique}.

Les extraits présentés dans cette annexe correspondent aux modules principaux du backend, du pipeline ML, du frontend et des tests. Chaque annexe indique le chemin relatif du fichier par rapport à la racine du projet.

% Config listings
\lstset{
  language=Python,
  basicstyle=\footnotesize\ttfamily,
  keywordstyle=\color{blue},
  commentstyle=\color{gray},
  stringstyle=\color{teal},
  numbers=left,
  numberstyle=\tiny,
  frame=single,
  breaklines=true,
  breakatwhitespace=false,
  tabsize=4,
  captionpos=b,
  aboveskip=1em,
  belowskip=1em,
}

% =============================================================================
\section{Annexe A1 -- Module ModbusManager}
\label{annexe:modbus}

\textbf{Fichier~:} \texttt{backend/modbus\_manager.py}

Ce module gère la communication Modbus RTU avec les six capteurs de température PT100 sur le bus RS485. Il implémente un verrou asynchrone pour l'exclusion mutuelle d'accès au bus, une reconnexion automatique après timeout, et un mappage des adresses esclaves vers les groupes et moules via \texttt{SENSOR\_MAP}.

\begin{lstlisting}
# modbus_manager.py
# Reads all temperature sensors via MODBUS RTU over RS485.
# Uses pymodbus with asyncio. Reads are strictly sequential because
# Modbus RTU is half-duplex: only one request/response exchange can be
# in flight at a time on the RS485 bus.

import asyncio
import logging
from dataclasses import dataclass
from datetime import datetime
from typing import Optional, List

from pymodbus.client import AsyncModbusSerialClient
from pymodbus.exceptions import ModbusException

import config

log = logging.getLogger(__name__)


@dataclass
class SensorReading:
    group_id    : int
    mold_id     : int
    position    : str
    temperature : Optional[float]
    status      : str        # "OK" | "ALERTE" | "CRITIQUE" | "ERREUR"
    threshold   : float
    deviation   : Optional[float]
    timestamp   : str


_client  : Optional[AsyncModbusSerialClient] = None
_bus_lock: Optional[asyncio.Lock]            = None   # one request at a time on RS485


async def init_modbus() -> bool:
    """Open the MODBUS serial connection. Returns True on success."""
    global _client, _bus_lock
    _bus_lock = asyncio.Lock()
    _client = AsyncModbusSerialClient(
        port     = config.MODBUS_PORT,
        baudrate = config.MODBUS_BAUDRATE,
        parity   = config.MODBUS_PARITY,
        stopbits = config.MODBUS_STOPBITS,
        bytesize = config.MODBUS_BYTESIZE,
        timeout  = config.MODBUS_TIMEOUT,
    )
    connected = await _client.connect()
    if connected:
        log.info("MODBUS connection established on %s", config.MODBUS_PORT)
    else:
        log.error("MODBUS connection failed on %s", config.MODBUS_PORT)
    return connected


async def _read_one(slave: int, register: int, retry: bool = True) -> Optional[float]:
    if _client is None or _bus_lock is None:
        return None
    for attempt in range(2 if retry else 1):
        try:
            async with _bus_lock:
                result = await _client.read_holding_registers(
                    register, count=1, slave=slave
                )
            if result.isError():
                log.warning("Sensor slave=%d reg=%d returned error", slave, register)
                continue
            raw = result.registers[0]
            return raw * config.TEMP_SCALE_FACTOR
        except asyncio.InvalidStateError:
            log.warning("Sensor slave=%d reg=%d InvalidStateError (attempt %d)", slave, register, attempt + 1)
        except asyncio.TimeoutError:
            log.warning("Sensor slave=%d reg=%d TimeoutError (attempt %d)", slave, register, attempt + 1)
        except (ModbusException, Exception) as exc:
            log.warning("Sensor slave=%d reg=%d %s (attempt %d)", slave, register, type(exc).__name__, attempt + 1)
    return None


async def read_all_sensors(calibration_temps: dict) -> List[SensorReading]:
    now      = datetime.now().isoformat(timespec='seconds')
    readings = []

    items = list(zip(config.SENSOR_MAP.keys(), config.SENSOR_MAP.values()))
    for idx, ((gid, mid), (slave, reg)) in enumerate(items):
        temp = await _read_one(slave, reg)
        offset = config.CALIBRATION_OFFSETS.get((gid, mid), 0)
        temp = temp + offset if temp is not None else None
        pos  = config.POSITION_MAP.get(mid, 'unknown')

        if idx < len(items) - 1:
            await asyncio.sleep(0.1)

        if temp is None:
            status    = 'ERREUR'
            deviation = None
        elif temp < config.T_MOLD_CRITICAL:
            status    = 'CRITIQUE'
            deviation = round(temp - config.T_HEATER, 3)
        elif temp < config.T_MOLD_WARNING:
            status    = 'ALERTE'
            deviation = round(temp - config.T_HEATER, 3)
        else:
            status    = 'OK'
            deviation = round(temp - config.T_HEATER, 3)

        readings.append(SensorReading(
            group_id    = gid,
            mold_id     = mid,
            position    = pos,
            temperature = round(temp, 2) if temp is not None else None,
            status      = status,
            threshold   = config.T_HEATER,
            deviation   = deviation,
            timestamp   = now,
        ))

    if readings and all(r.temperature is None for r in readings):
        now_ts = datetime.now().timestamp()
        elapsed = now_ts - getattr(read_all_sensors, '_last_reconnect', 0)
        if elapsed > 30:
            read_all_sensors._last_reconnect = now_ts
            log.warning("All sensors returned None -- attempting Modbus reconnection")
            await close_modbus()
            await init_modbus()
        else:
            log.debug("All sensors None, reconnect debounced (%.0fs since last)", elapsed)

    return readings


async def read_heater_temp() -> Optional[float]:
    if config.HEATER_TEMP_SENSOR is None:
        return None
    slave, reg = config.HEATER_TEMP_SENSOR
    return await _read_one(slave, reg)


async def close_modbus():
    global _client
    if _client:
        _client.close()
        _client = None
        log.info("MODBUS connection closed")
\end{lstlisting}

% =============================================================================
\section{Annexe A2 -- Modèle Grey-Box}
\label{annexe:greybox}

\textbf{Fichier~:} \texttt{ml/grey\_box.py}

Le modèle Grey-Box implémente le \textit{soft-sensing} de l'épaisseur de calcaire à partir des mesures de température et de débit. Il calcule la résistance thermique, l'épaisseur, la dégradation relative et émet un niveau d'urgence.

\begin{lstlisting}
import math
import logging
from typing import Optional, Dict, Tuple

import config

log = logging.getLogger(__name__)

PIPE_AREA = math.pi * config.PIPE_LENGTH * config.PIPE_DIAMETER


class GreyBoxModel:

    def __init__(self):
        self.calibration_temps: Dict[Tuple, float] = {}
        self.delta_T_normal:    Dict[Tuple, float] = {}

    def set_calibration(self, group_id: int, mold_id: int, T_mold_jour1: float):
        key = (group_id, mold_id)
        self.calibration_temps[key] = T_mold_jour1
        self.delta_T_normal[key]    = round(config.T_HEATER - T_mold_jour1, 4)
        log.debug(
            "Calibration set mold (%d,%d): T_jour1=%.2f  delta_T_normal=%.4f",
            group_id, mold_id, T_mold_jour1, self.delta_T_normal[key]
        )

    def compute(self, group_id: int, mold_id: int, T_mold: float, flow_lpm: float) -> Dict:
        key = (group_id, mold_id)

        flow_m3s    = (flow_lpm / 60.0) / 1000.0
        flow_mold   = flow_m3s

        delta_T_measured  = config.T_HEATER - T_mold
        delta_T_norm      = self.delta_T_normal.get(key, 1.5)

        delta_T_calcaire  = max(0.0, delta_T_measured - delta_T_norm)

        Q = flow_mold * config.RHO_WATER * config.CP_WATER * delta_T_measured
        Q = max(Q, 1e-6)

        R_calcaire = delta_T_calcaire / Q
        epaisseur_mm = R_calcaire * config.LAMBDA_CALCAIRE * PIPE_AREA * 1000.0

        delta_T_max   = (config.T_HEATER - config.T_MOLD_CRITICAL) - delta_T_norm
        delta_T_max   = max(delta_T_max, 0.1)
        degradation_pct = min((delta_T_calcaire / delta_T_max) * 100.0, 100.0)

        urgence = self._classify_urgency(T_mold)

        return {
            'delta_T_measured':  round(delta_T_measured,  4),
            'delta_T_calcaire':  round(delta_T_calcaire,  4),
            'Q':                 round(Q,                  2),
            'R_calcaire':        round(R_calcaire,         6),
            'epaisseur_mm':      round(epaisseur_mm,       4),
            'urgence':           urgence,
            'degradation_pct':   round(degradation_pct,    1),
        }

    @staticmethod
    def _classify_urgency(T_mold: float) -> str:
        if T_mold >= 42.0:
            return 'OK'
        elif T_mold >= 41.5:
            return 'FAIBLE'
        elif T_mold >= 41.0:
            return 'MOYEN'
        elif T_mold >= 40.5:
            return 'HAUTE'
        else:
            return 'URGENT'
\end{lstlisting}

% =============================================================================
\section{Annexe A3 -- Détection d'anomalies par Isolation Forest}
\label{annexe:anomaly}

\textbf{Fichier~:} \texttt{ml/anomaly\_detector.py}

Le détecteur d'anomalies utilise \texttt{IsolationForest} de scikit-learn avec 200 estimateurs et un taux de contamination de 5\%. L'extraction des huit caractéristiques s'effectue sur une fenêtre glissante de 30 secondes.

\begin{lstlisting}
import logging
import pickle
import os
import numpy as np
from typing import Optional, List, Dict, Tuple

from sklearn.ensemble import IsolationForest
from sklearn.preprocessing import StandardScaler

import config

log = logging.getLogger(__name__)

_ML_DIR = os.path.dirname(os.path.abspath(__file__))
_PROJECT_ROOT = os.path.dirname(_ML_DIR)
MODEL_PATH  = os.path.join(_PROJECT_ROOT, 'models', 'isolation_forest.pkl')
SCALER_PATH = os.path.join(_PROJECT_ROOT, 'models', 'scaler_if.pkl')


class AnomalyDetector:

    FEATURE_NAMES = [
        'slope_T_mold', 'variance_T_mold', 'affected_molds_ratio',
        'sudden_drop_flag', 'flow_rate', 'flow_variance',
        'delta_T_calcaire_mean', 'autocorr_lag1',
    ]

    def __init__(self):
        self.model:   Optional[IsolationForest] = None
        self.scaler:  Optional[StandardScaler]  = None
        self.trained: bool                      = False
        os.makedirs('models', exist_ok=True)
        self._load()

    def extract_features(self, temp_history, flow_history, delta_T_calcaires) -> Optional[np.ndarray]:
        window  = config.FEATURE_WINDOW_SECONDS
        all_temps = []
        slopes, variances, autocorrs = [], [], []
        affected = 0

        for key, hist in temp_history.items():
            if len(hist) < 10:
                continue
            arr = np.array(hist[-window:], dtype=float)
            all_temps.extend(arr.tolist())
            x  = np.arange(len(arr))
            slope = np.polyfit(x, arr, 1)[0] if len(arr) > 1 else 0.0
            slopes.append(slope)
            variances.append(float(np.var(arr)))
            if len(arr) > 2:
                with np.errstate(invalid='ignore', divide='ignore'):
                    ac = float(np.corrcoef(arr[:-1], arr[1:])[0, 1])
                autocorrs.append(ac if not np.isnan(ac) else 0.0)
            if arr[-1] < config.T_MOLD_CRITICAL:
                affected += 1

        if not slopes:
            return None

        sudden_drop = 0.0
        for key, hist in temp_history.items():
            if len(hist) >= 120:
                if hist[-1] - hist[-120] < -1.0:
                    sudden_drop = 1.0
                    break

        n_molds = max(len(temp_history), 1)
        flow_means = [float(np.mean(v[-window:])) for v in flow_history.values() if v] or [config.FLOW_DEFAULT_LPM]
        flow_vars  = [float(np.var(v[-window:]))  for v in flow_history.values() if v] or [0.0]
        flow_mean  = float(np.mean(flow_means))
        flow_var   = float(np.mean(flow_vars))
        dT_mean = float(np.mean(list(delta_T_calcaires.values()))) if delta_T_calcaires else 0.0

        features = np.array([[
            float(np.mean(slopes)), float(np.mean(variances)),
            affected / n_molds, sudden_drop,
            flow_mean, flow_var, dT_mean,
            float(np.mean(autocorrs)) if autocorrs else 0.0,
        ]])
        return features

    def train(self, feature_matrix: np.ndarray):
        self.scaler = StandardScaler()
        X_scaled    = self.scaler.fit_transform(feature_matrix)
        self.model = IsolationForest(
            n_estimators=200, contamination=0.05, random_state=42, n_jobs=-1,
        )
        self.model.fit(X_scaled)
        self.trained = True
        self._save()
        log.info("Isolation Forest trained on %d samples", len(feature_matrix))

    def predict(self, features: np.ndarray) -> Dict:
        if not self.trained or self.model is None or self.scaler is None:
            return {'anomaly_detected': False, 'anomaly_score': None}
        X_scaled = self.scaler.transform(features)
        label    = self.model.predict(X_scaled)[0]
        score    = self.model.score_samples(X_scaled)[0]
        return {
            'anomaly_detected': bool(label == -1),
            'anomaly_score':    float(round(score, 4)),
        }

    def _save(self):
        try:
            with open(MODEL_PATH,  'wb') as f: pickle.dump(self.model,  f)
            with open(SCALER_PATH, 'wb') as f: pickle.dump(self.scaler, f)
            log.info("Isolation Forest model saved")
        except Exception as exc:
            log.error("Could not save Isolation Forest: %s", exc)

    def _load(self):
        try:
            if os.path.exists(MODEL_PATH) and os.path.exists(SCALER_PATH):
                with open(MODEL_PATH,  'rb') as f: self.model  = pickle.load(f)
                with open(SCALER_PATH, 'rb') as f: self.scaler = pickle.load(f)
                self.trained = True
                log.info("Isolation Forest model loaded from disk")
        except Exception as exc:
            log.warning("Could not load Isolation Forest: %s", exc)
\end{lstlisting}

% =============================================================================
\section{Annexe A4 -- Classifieur hybride N1/N2}
\label{annexe:classifier}

\textbf{Fichier~:} \texttt{ml/cause\_classifier.py}

Le classifieur de causes implémente l'architecture à deux niveaux~: règles physiques déterministes (N1) puis Random Forest (N2) pour les cas ambigus. Il inclut la méthode \texttt{auto\_label()} pour l'étiquetage automatique.

\begin{lstlisting}
import logging
import pickle
import os
import numpy as np
from typing import Optional, Dict, Tuple, List

from sklearn.ensemble import RandomForestClassifier
from sklearn.preprocessing import LabelEncoder

import config

log = logging.getLogger(__name__)

_ML_DIR = os.path.dirname(os.path.abspath(__file__))
_PROJECT_ROOT = os.path.dirname(_ML_DIR)
MODEL_PATH   = os.path.join(_PROJECT_ROOT, 'models', 'random_forest.pkl')
ENCODER_PATH = os.path.join(_PROJECT_ROOT, 'models', 'label_encoder.pkl')

CLASSES = [
    'CALCAIRE_TUYAUX', 'HEATER_POMPE_HS', 'HEATER_RESISTANCE_HS',
    'NIVEAU_BAS_VANNE_PANNE', 'BULLES_AIR', 'FUITE_CIRCUIT', 'ISOLATION_DEGRADEE'
]


class CauseClassifier:

    FEATURE_NAMES = [
        'slope_T_mold', 'variance_T_mold', 'affected_molds_ratio',
        'sudden_drop_flag', 'flow_rate', 'flow_drop_flag', 'flow_variance',
        'delta_T_calcaire_slope', 'drift_R_squared', 'autocorr_lag1',
    ]

    def __init__(self):
        self.model:   Optional[RandomForestClassifier] = None
        self.encoder: Optional[LabelEncoder]           = None
        self.trained: bool                             = False
        os.makedirs('models', exist_ok=True)
        self._load()

    @staticmethod
    def physical_rules(affected_ratio, sudden_drop, flow_rate, flow_drop,
                       temp_heater=45.0, nominal_flow=16.5) -> Optional[Dict]:
        if temp_heater < 44.0 and affected_ratio > 0.8:
            return {'cause': 'HEATER_RESISTANCE_HS', 'confidence': 1.0, 'method': 'physical_rule'}
        if affected_ratio > 0.8 and flow_drop:
            return {'cause': 'HEATER_POMPE_HS', 'confidence': 1.0, 'method': 'physical_rule'}
        if affected_ratio > 0.7 and flow_rate < 0.5:
            return {'cause': 'NIVEAU_BAS_VANNE_PANNE', 'confidence': 1.0, 'method': 'physical_rule'}
        return None

    def predict(self, features: np.ndarray) -> Dict:
        if not self.trained or self.model is None:
            return {'cause': 'CAUSE_INDETERMINEE', 'confidence': 0.0, 'method': 'default', 'proba_dict': {}}
        proba  = self.model.predict_proba(features)[0]
        idx    = int(np.argmax(proba))
        cause  = self.encoder.inverse_transform([idx])[0]
        max_prob = float(round(float(proba[idx]), 3))
        proba_dict = {cls: float(round(p, 3)) for cls, p in zip(self.encoder.classes_, proba)}
        if max_prob < config.RF_CONFIDENCE_THRESHOLD:
            return {'cause': 'CAUSE_INDETERMINEE', 'confidence': max_prob, 'proba_dict': proba_dict, 'method': 'low_confidence'}
        return {'cause': cause, 'confidence': max_prob, 'proba_dict': proba_dict, 'method': 'random_forest'}

    def train(self, X: np.ndarray, y: List[str]):
        self.encoder = LabelEncoder()
        y_encoded = self.encoder.fit_transform(y)
        self.model = RandomForestClassifier(
            n_estimators=100, class_weight='balanced', max_depth=10, random_state=42, n_jobs=-1,
        )
        self.model.fit(X, y_encoded)
        self.trained = True
        self._save()
        log.info("Random Forest trained on %d samples, classes: %s", len(X), list(self.encoder.classes_))

    def feature_importances(self) -> Dict[str, float]:
        if self.model is None:
            return {}
        return dict(zip(self.FEATURE_NAMES, self.model.feature_importances_))

    @staticmethod
    def auto_label(affected_ratio, sudden_drop, flow_drop, flow_rate,
                   variance, R_squared, delta_T_calcaire_slope,
                   temp_heater=45.0, variance_threshold=0.1, nominal_flow=16.5) -> str:
        if temp_heater < 44.0 and affected_ratio > 0.8:
            return 'HEATER_RESISTANCE_HS'
        if affected_ratio > 0.8 and flow_drop:
            return 'HEATER_POMPE_HS'
        if affected_ratio > 0.7 and flow_rate < 0.5:
            return 'NIVEAU_BAS_VANNE_PANNE'
        if variance > variance_threshold and R_squared < 0.3 and affected_ratio < 0.4:
            return 'BULLES_AIR'
        if delta_T_calcaire_slope > 0.03 and R_squared > 0.85:
            return 'CALCAIRE_TUYAUX'
        if affected_ratio < 0.3 and R_squared > 0.7:
            return 'ISOLATION_DEGRADEE'
        return 'NORMAL'

    def _save(self):
        try:
            with open(MODEL_PATH,   'wb') as f: pickle.dump(self.model,   f)
            with open(ENCODER_PATH, 'wb') as f: pickle.dump(self.encoder, f)
            log.info("Random Forest saved")
        except Exception as exc:
            log.error("Could not save Random Forest: %s", exc)

    def _load(self):
        try:
            if os.path.exists(MODEL_PATH) and os.path.exists(ENCODER_PATH):
                with open(MODEL_PATH,   'rb') as f: self.model   = pickle.load(f)
                with open(ENCODER_PATH, 'rb') as f: self.encoder = pickle.load(f)
                self.trained = True
                log.info("Random Forest loaded from disk")
        except Exception as exc:
            log.warning("Could not load Random Forest: %s", exc)
\end{lstlisting}

% =============================================================================
\section{Annexe A5 -- Prédiction Ridge Regression}
\label{annexe:ridge}

\textbf{Fichier~:} \texttt{ml/ridge\_predictor.py}

Le module \texttt{RidgePredictor} implémente la régression Ridge polynomiale avec validation croisée pour la prédiction de durée de vie résiduelle (TTF). L'intervalle de confiance à 90\% est estimé par bootstrap (1000 échantillons).

\begin{lstlisting}
import logging
import pickle
import os
import numpy as np
from datetime import datetime, timedelta
from typing import Optional, Dict, List, Tuple

from sklearn.linear_model import Ridge
from sklearn.preprocessing import PolynomialFeatures

import config

log = logging.getLogger(__name__)

_ML_DIR = os.path.dirname(os.path.abspath(__file__))
_PROJECT_ROOT = os.path.dirname(_ML_DIR)
MODEL_DIR = os.path.join(_PROJECT_ROOT, 'models', 'ridge')


class RidgePredictor:

    def __init__(self, group_id: int, mold_id: int, delta_T_max: float):
        self.group_id   = group_id
        self.mold_id    = mold_id
        self.delta_T_max = max(delta_T_max, 0.1)
        self.poly        = PolynomialFeatures(degree=2, include_bias=True)
        self.model:  Optional[Ridge] = None
        self.X_data: Optional[np.ndarray] = None
        self.y_data: Optional[np.ndarray] = None
        self.n_days: int = 0
        os.makedirs(MODEL_DIR, exist_ok=True)
        self._load()

    def fit(self, daily_records: List[Dict]):
        if len(daily_records) < config.RIDGE_MIN_DAYS:
            log.info("Ridge mold (%d,%d): not enough data (%d < %d days)",
                     self.group_id, self.mold_id, len(daily_records), config.RIDGE_MIN_DAYS)
            return
        X_raw = np.array([r['day_offset'] for r in daily_records]).reshape(-1, 1)
        y     = np.array([r['value']      for r in daily_records])
        X_poly    = self.poly.fit_transform(X_raw)
        self.model = Ridge(alpha=1.0)
        self.model.fit(X_poly, y)
        self.X_data = X_raw
        self.y_data = y
        self.n_days = len(daily_records)
        self._save()
        log.info("Ridge mold (%d,%d) trained on %d days", self.group_id, self.mold_id, self.n_days)

    def predict_maintenance(self) -> Optional[Dict]:
        if self.model is None or self.X_data is None:
            return None
        N     = len(self.X_data)
        start = int(self.X_data[-1]) + 1
        X_future = np.arange(start, start + 366).reshape(-1, 1)
        X_f_poly = self.poly.transform(X_future)
        y_pred = self.model.predict(X_f_poly)
        idx    = self._find_crossing(y_pred)
        if idx is None:
            return None
        central = int(idx)
        crossing_days = []
        for _ in range(config.BOOTSTRAP_N):
            indices  = np.random.choice(N, size=N, replace=True)
            X_boot   = self.poly.transform(self.X_data[indices])
            y_boot   = self.y_data[indices]
            m        = Ridge(alpha=1.0)
            m.fit(X_boot, y_boot)
            y_b      = m.predict(X_f_poly)
            i_b      = self._find_crossing(y_b)
            if i_b is not None:
                crossing_days.append(int(i_b))
        if not crossing_days:
            return None
        borne_basse = int(np.percentile(crossing_days, 5))
        borne_haute = int(np.percentile(crossing_days, 95))
        median      = int(np.percentile(crossing_days, 50))
        predicted_date = (datetime.today() + timedelta(days=borne_basse)).strftime('%d/%m/%Y')
        return {
            'jours_maintenance': median, 'borne_basse': borne_basse,
            'borne_haute': borne_haute, 'predicted_date': predicted_date,
            'n_bootstrap': len(crossing_days),
        }

    def _find_crossing(self, y_pred: np.ndarray) -> Optional[int]:
        crossings = np.where(y_pred >= self.delta_T_max)[0]
        return int(crossings[0]) if len(crossings) > 0 else None

    def _save(self):
        path = os.path.join(MODEL_DIR, f'ridge_{self.group_id}_{self.mold_id}.pkl')
        try:
            with open(path, 'wb') as f:
                pickle.dump({'model': self.model, 'X': self.X_data, 'y': self.y_data,
                             'n': self.n_days, 'poly': self.poly}, f)
        except Exception as exc:
            log.error("Could not save Ridge mold (%d,%d): %s", self.group_id, self.mold_id, exc)

    def _load(self):
        path = os.path.join(MODEL_DIR, f'ridge_{self.group_id}_{self.mold_id}.pkl')
        try:
            if os.path.exists(path):
                with open(path, 'rb') as f:
                    d = pickle.load(f)
                self.model  = d['model']
                self.X_data = d['X']
                self.y_data = d['y']
                self.n_days = d['n']
                self.poly   = d['poly']
                log.info("Ridge mold (%d,%d) loaded from disk", self.group_id, self.mold_id)
        except Exception as exc:
            log.warning("Could not load Ridge mold (%d,%d): %s", self.group_id, self.mold_id, exc)
\end{lstlisting}

% =============================================================================
\section{Annexe A6 -- Système d'alertes}
\label{annexe:alerting}

\textbf{Fichier~:} \texttt{backend/alerting.py}

Le module d'alertes envoie des notifications par deux canaux~: Telegram (API HTTP) pour les alertes WARNING, et Email (SMTP Gmail) pour les alertes CRITICAL et SYSTEM.

\begin{lstlisting}
import logging
import smtplib
from email.message import EmailMessage
from typing import List

import requests

import config

log = logging.getLogger(__name__)

TELEGRAM_API = f"https://api.telegram.org/bot{config.TELEGRAM_BOT_TOKEN}/sendMessage"


def send_telegram(chat_id: str, message: str) -> bool:
    try:
        resp = requests.post(
            TELEGRAM_API,
            json={"chat_id": chat_id, "text": message, "parse_mode": "Markdown"},
            timeout=10,
        )
        if resp.status_code != 200:
            log.warning("Telegram API error (%d): %s", resp.status_code, resp.text)
            return False
        return True
    except Exception as exc:
        log.warning("Telegram send failed: %s", exc)
        return False


def send_email(subject: str, body: str) -> bool:
    try:
        msg = EmailMessage()
        msg.set_content(body)
        msg["Subject"] = subject
        msg["From"] = config.EMAIL_FROM
        msg["To"] = config.EMAIL_TO
        msg["Cc"] = config.EMAIL_FROM
        with smtplib.SMTP(config.SMTP_HOST, config.SMTP_PORT, timeout=15) as server:
            server.starttls()
            server.login(config.SMTP_USER, config.SMTP_PASSWORD)
            server.send_message(msg)
        return True
    except Exception as exc:
        log.warning("Email send failed: %s", exc)
        return False


def _format_mold_line(m: dict) -> str:
    idx = m.get('global_idx', m.get('mold_id', '?'))
    pos = config.POSITION_MAP.get(m.get('mold_id'), '?')
    nomenclature = m.get('nomenclature', '') or config.MOLD_NOMENCLATURE.get(
        (m.get('group_id'), m.get('mold_id')), '')
    temp = m.get('temperature')
    temp_str = f"{temp:.1f} °C" if temp is not None else "-- °C"
    status = m.get('status', 'ERREUR')
    name_part = f" [{nomenclature}]" if nomenclature else ""
    return f"   Moule {idx}{name_part} ({pos}) : {temp_str} [{status}]"


def send_alert(severity: str, affected_molds: List[int], mold_readings: List[dict]) -> None:
    ts = __import__("datetime").datetime.now().strftime("%d/%m/%Y %H:%M")
    header = f"{'🔴 *ALERTE CRITIQUE*' if severity == 'CRITICAL' else '⚠️ *Alerte Température*'}"
    molds_str = ', '.join(str(m) for m in affected_molds) if affected_molds else '-'
    temps_lines = '\n'.join(_format_mold_line(m) for m in mold_readings)
    message = (
        f"{header}\n\n"
        f"🕐 {ts}\n"
        f"🔸 Moules affectés : {molds_str}\n"
        + (f"\n🔸 Températures :\n{temps_lines}" if temps_lines else "")
    )
    if severity == "WARNING":
        send_telegram(config.TELEGRAM_OPERATORS_ID, message)
    elif severity == "CRITICAL":
        send_telegram(config.TELEGRAM_OPERATORS_ID, message)
        send_telegram(config.TELEGRAM_CHEF_ID, message)
        email_body = (
            f"Alerte Critique - Supervision Thermique\n\n"
            f"Date : {ts}\n"
            f"Moules affectés : {molds_str}\n"
            + (f"\nTempératures :\n{temps_lines}" if temps_lines else "")
        )
        send_email(f"ALERTE - Supervision Thermique - {ts}", email_body)
\end{lstlisting}

% =============================================================================
\section{Annexe A7 -- Acquisition du débitmètre YF-S201}
\label{annexe:flow}

\textbf{Fichier~:} \texttt{backend/flow\_sensor.py}

Le module \texttt{FlowSensor} utilise une interruption matérielle GPIO via \texttt{gpiozero} pour compter les impulsions du débitmètre YF-S201. La conversion fréquence-débit suit \(Q = F / 7,5\).

\begin{lstlisting}
import time
import logging
import threading

try:
    import gpiozero
    HAS_GPIO = True
except ImportError:
    HAS_GPIO = False
    logging.warning("gpiozero not installed. Flow sensor will use default value.")

import config

log = logging.getLogger(__name__)


class FlowSensor:
    def __init__(self, pin, k_factor=7.5):
        self.pin = pin
        self.k_factor = k_factor
        self.pulse_count = 0
        self.lock = threading.Lock()
        self.sensor = None
        self.last_read_time = time.time()
        self._init_sensor()

    def _init_sensor(self):
        if not HAS_GPIO:
            log.warning("gpiozero not available, flow sensor disabled")
            return
        try:
            self.sensor = gpiozero.DigitalInputDevice(self.pin, bounce_time=0.01)
            self.sensor.when_activated = self._pulse
            log.info(f"Flow sensor initialized on GPIO {self.pin}")
        except Exception as e:
            log.error(f"Failed to initialize flow sensor on GPIO {self.pin}: {e}")
            self.sensor = None

    def _pulse(self):
        with self.lock:
            self.pulse_count += 1

    def read_lpm(self) -> float:
        if not self.sensor:
            return config.FLOW_DEFAULT_LPM
        now = time.time()
        elapsed = now - self.last_read_time
        self.last_read_time = now
        with self.lock:
            pulses = self.pulse_count
            self.pulse_count = 0
        frequency = pulses / elapsed if elapsed > 0 else 0
        flow_rate = frequency / self.k_factor
        if flow_rate > 30:
            log.warning(f"Flow rate spike detected: {flow_rate:.2f} L/min")
            flow_rate = config.FLOW_DEFAULT_LPM
        return flow_rate

    def close(self):
        if self.sensor:
            self.sensor.close()
            log.info("Flow sensor closed")
\end{lstlisting}

% =============================================================================
\section{Annexe A8 -- Connexion WebSocket (Frontend)}
\label{annexe:websocket}

\textbf{Fichier~:} \texttt{frontend/src/hooks/useWebSocket.js}

Ce hook React gère la connexion WebSocket bidirectionnelle avec le backend FastAPI. Il assure la reconnexion automatique avec \textit{exponential backoff} et distribue les données mises à jour aux trois onglets.

\begin{lstlisting}[language=JavaScript]
// useWebSocket.js
// Manages the WebSocket connection to the FastAPI backend.
// Handles automatic reconnection with exponential backoff.

import { useState, useEffect, useRef, useCallback } from 'react'

const WS_URL = '/ws'

export function useWebSocket() {
  const [data, setData]           = useState(null)
  const [connected, setConnected] = useState(false)
  const wsRef                     = useRef(null)
  const retryDelay                = useRef(1000)
  const retryTimer                = useRef(null)
  const reconnectRef = useRef(true)

  const connect = useCallback(() => {
    if (wsRef.current && wsRef.current.readyState === WebSocket.OPEN) return
    const ws = new WebSocket(WS_URL)
    wsRef.current = ws

    ws.onopen = () => {
      setConnected(true)
      retryDelay.current = 1000
    }

    ws.onmessage = (event) => {
      try {
        const parsed = JSON.parse(event.data)
        if (parsed && parsed.t === 'ping') return
        setData(parsed)
      } catch (err) {
        console.error('WebSocket parse error:', err)
      }
    }

    ws.onclose = () => {
      setConnected(false)
      if (!reconnectRef.current) return
      retryTimer.current = setTimeout(() => {
        retryDelay.current = Math.min(retryDelay.current * 2, 30000)
        if (reconnectRef.current) connect()
      }, retryDelay.current)
    }

    ws.onerror = () => { ws.close() }
  }, [])

  useEffect(() => {
    reconnectRef.current = true
    connect()
    return () => {
      reconnectRef.current = false
      clearTimeout(retryTimer.current)
      if (wsRef.current) {
        wsRef.current.onclose = null
        wsRef.current.close()
      }
    }
  }, [connect])

  return { data, connected }
}
\end{lstlisting}

% =============================================================================
\section{Annexe A9 -- Démonstrateur CLI (demo\_cli)}
\label{annexe:demo_cli}

\textbf{Fichier~:} \texttt{tests/demo\_cli.py}

L'outil \texttt{demo\_cli.py} est un démonstrateur en ligne de commande qui se connecte à l'API FastAPI du backend et change dynamiquement le mode de simulation. Il propose six scénarios de défaillance ainsi qu'une auto-séquence.

\begin{lstlisting}
import sys
import os
sys.path.insert(0, os.path.join(os.path.dirname(__file__), '..', 'backend'))

import requests

BASE_URL = "http://localhost:8001"

MODES = [
    ('1', 'NORMAL',       'all temperatures stable around 44-45°C'),
    ('2', 'GRADUAL_DROP', 'slow temperature decay (-0.03°C per call)'),
    ('3', 'SUDDEN_DROP',  'sudden -2.5°C drop at 60th call'),
    ('4', 'HEATER_FAIL',  'all molds dropping below 42°C (heater failure)'),
    ('5', 'PUMP_FAIL',    'global drop + erratic readings (pump failure)'),
    ('6', 'NOISY',        'normal with 10% random ERREUR readings'),
]


def set_mode(mode: str):
    try:
        resp = requests.post(f"{BASE_URL}/api/sim/mode", json={"mode": mode}, timeout=3)
        if resp.status_code == 200:
            print(f"  \u2713 Mode chang\u00e9 : {mode}")
        else:
            data = resp.json()
            print(f"  \u2717 Erreur {resp.status_code} : {data.get('detail', resp.text)}")
    except requests.ConnectionError:
        print(f"  \u2717 Impossible de se connecter au backend ({BASE_URL})")
    except Exception as e:
        print(f"  \u2717 Erreur : {e}")


def auto_sequence():
    modes_seq = [
        ('NORMAL', 2), ('GRADUAL_DROP', 2), ('SUDDEN_DROP', 1),
        ('HEATER_FAIL', 1), ('NORMAL', 0.5), ('PUMP_FAIL', 1),
        ('NOISY', 1), ('NORMAL', 0),
    ]
    print("\n  S\u00e9quence automatique de d\u00e9monstration :")
    for mode, pause in modes_seq:
        set_mode(mode)
        if pause > 0:
            import time
            time.sleep(pause)
    print("  Termin\u00e9.\n")


def clear_screen():
    os.system('cls' if os.name == 'nt' else 'clear')


def main():
    while True:
        clear_screen()
        print("=" * 60)
        print("  SUPERVISION THERMIQUE \u2014 Contr\u00f4le de la Simulation")
        print("=" * 60)
        print()
        for key, mode, desc in MODES:
            print(f"  {key}. {mode:<15s}   {desc}")
        print()
        print("  A. S\u00e9quence automatique (encha\u00eene tous les modes)")
        print("  Q. Quitter")
        print()
        choice = input("  Choix : ").strip().upper()

        if choice == 'Q':
            print("  Au revoir.")
            break
        if choice == 'A':
            auto_sequence()
            input("  Appuyez sur Entr\u00e9e pour continuer...")
            continue

        matched = [m for k, m, d in MODES if k == choice]
        if matched:
            set_mode(matched[0])
            input("  Appuyez sur Entr\u00e9e pour continuer...")
        else:
            print("  Choix invalide.")
            input("  Appuyez sur Entr\u00e9e...")
    print()


if __name__ == '__main__':
    main()
\end{lstlisting}

% =============================================================================
\section{Annexe A10 -- Configuration générale}
\label{annexe:config}

\textbf{Fichier~:} \texttt{backend/config.py}

Ce fichier centralise l'ensemble des paramètres de configuration~: adresses Modbus, mappage des esclaves, paramètres InfluxDB, seuils de température, paramètres ML, configuration SMTP et jeton Telegram.

\begin{lstlisting}
# config.py
# Configuration pour le backend

import os

# -- MODBUS -------------------------------------------------------------------
MODBUS_PORT      = '/dev/ttyUSB0'
MODBUS_BAUDRATE  = 9600
MODBUS_PARITY    = 'N'
MODBUS_STOPBITS  = 1
MODBUS_BYTESIZE  = 8
MODBUS_TIMEOUT   = 1

SENSOR_MAP = {
    (1, 1): (1, 0),   # Heater 1 -- moule 1 (gauche)
    (2, 1): (2, 0),   # Heater 2 -- moule 1 (gauche)
    (2, 2): (3, 0),   # Heater 2 -- moule 2 (centre)
    (3, 1): (4, 0),   # Heater 3 -- moule 1 (gauche)
    (3, 2): (5, 0),   # Heater 3 -- moule 2 (centre)
    (4, 1): (7, 0),   # Heater 4 -- moule 1 (gauche)
}

POSITION_MAP = { 1: 'gauche', 2: 'droite' }
TEMP_SCALE_FACTOR = 0.1

MOLD_NOMENCLATURE = {
    (1, 1): 'RBA_178 #03', (2, 1): 'RBA_178 #01', (2, 2): 'RBA_178 #02',
    (3, 1): 'RBA_174 #02', (3, 2): 'RBA_174 #06', (4, 1): 'RBA_174 #08',
}

# -- TEMPERATURE THRESHOLDS ---------------------------------------------------
T_HEATER        = 45.0
T_TOLERANCE     = 3.0
T_MOLD_WARNING  = 42.0
T_MOLD_CRITICAL = 40.0

# -- INFLUXDB -----------------------------------------------------------------
INFLUX_URL    = 'http://localhost:8086'
INFLUX_TOKEN  = '4tAf1cnPev0dM0rUUjTWpcVbMWZviB9eQeX04-W0P2Gj_Xu7JLCbvhPjMXEAfD4GkKgrKA6Y9c7otYyFcrvmQQ=='
INFLUX_ORG    = 'temperature-monitoring'
INFLUX_BUCKET = 'sensors'

# -- ALERTES ------------------------------------------------------------------
TELEGRAM_BOT_TOKEN    = os.getenv('TELEGRAM_BOT_TOKEN', 'ton_token_ici')
TELEGRAM_OPERATORS_ID = os.getenv('TELEGRAM_OPERATORS_ID', '-100...')
TELEGRAM_CHEF_ID      = os.getenv('TELEGRAM_CHEF_ID', '-100...')

SMTP_HOST     = os.getenv('SMTP_HOST', 'smtp.gmail.com')
SMTP_PORT     = int(os.getenv('SMTP_PORT', '587'))
SMTP_USER     = os.getenv('SMTP_USER', 'ton.email@gmail.com')
SMTP_PASSWORD = os.getenv('SMTP_PASSWORD', 'mdp_application')
EMAIL_FROM    = os.getenv('EMAIL_FROM', 'ton.email@gmail.com')
EMAIL_TO      = os.getenv('EMAIL_TO', 'chef.equipe@yazaki.ma')

# -- FLOW SENSOR --------------------------------------------------------------
FLOW_SENSOR_PINS = { 1: 17 }
FLOW_DEFAULT_LPM     = 16.5
YF_S201_PULSE_FACTOR = 7.5

# -- GREY-BOX -----------------------------------------------------------------
PIPE_DIAMETER   = 0.013
PIPE_LENGTH     = 3.0
LAMBDA_CALCAIRE = 1.0
N_MOLDS         = 6
RHO_WATER       = 1000.0
CP_WATER        = 4186.0

# -- ML -----------------------------------------------------------------------
FEATURE_WINDOW_SECONDS = 30
RIDGE_MIN_DAYS         = 7
BOOTSTRAP_N            = 1000
RETRAIN_HOUR           = 2

# -- WEBSOCKET ----------------------------------------------------------------
WS_HOST = '0.0.0.0'
WS_PORT = 8001

# -- ACQUISITION --------------------------------------------------------------
ACQUISITION_HZ = 1

# -- AMDEC --------------------------------------------------------------------
AMDEC_FAILURE_MODES = {
    'NIVEAU_BAS_VANNE_PANNE': {
        'criticite': 180, 'priorite': 1,
        'actions': [
            "Verifier niveau reservoir visuellement",
            "Tester vanne appoint manuellement",
            "Verifier si vanne s'ouvre quand niveau descend",
            "Verifier commande electrique vanne (automate)",
            "Si vanne bloquee : Debloquer ou remplacer",
            "Remplir reservoir manuellement si necessaire"
        ]
    },
    'HEATER_RESISTANCE_HS': {
        'criticite': 160, 'priorite': 2,
        'actions': [
            "Mesurer T_heater (doit etre 45C +/- 1C)",
            "Si T_heater < 44C : resistance probablement HS",
            "Verifier alimentation electrique resistance",
            "Tester continuite resistance (multimetre)",
            "Verifier regulation temperature PID",
            "Remplacement resistance urgent"
        ]
    },
    'CALCAIRE_TUYAUX': {
        'criticite': 96, 'priorite': 3,
        'actions': [
            "Mesurer debit d'eau actuel avec debitmetre",
            "Calculer epaisseur calcaire (modele Grey-Box)",
            "Si epaisseur > 1.5mm : Planifier detartrage chimique sous 7 jours",
            "Surveiller evolution debit toutes les 4h"
        ]
    },
    'HEATER_POMPE_HS': {
        'criticite': 90, 'priorite': 4,
        'actions': [
            "Verifier debit pompe heater (nominal: 16.5 L/min)",
            "Ecouter bruits anormaux/vibrations pompe",
            "Verifier alimentation electrique pompe",
            "Inspecter roue pompe (usure, blocage, cavitation)",
            "Si pompe HS : Arret production, remplacement urgent"
        ]
    },
    'BULLES_AIR': {
        'criticite': 90, 'priorite': 5,
        'actions': [
            "Purger circuit hydraulique immediatement",
            "Verifier niveau reservoir (doit etre > 50%)",
            "Inspecter joints pompe (entrees air possibles)",
            "Controler vannes d'admission",
            "Surveiller stabilite temperature 30 min post-purge"
        ]
    },
    'FUITE_CIRCUIT': {
        'criticite': 60, 'priorite': 6,
        'actions': [
            "Inspecter visuellement circuit complet",
            "Rechercher fuites (joints, raccords, tuyaux)",
            "Mesurer pression circuit si possible",
            "Reparer fuite identifiee",
            "Tester etancheite apres reparation"
        ]
    },
    'ISOLATION_DEGRADEE': {
        'criticite': 45, 'priorite': 7,
        'actions': [
            "Inspecter isolation moule visuellement"
        ]
    }
}

AMDEC_CAUSE_MAPPING = {
    'POMPE_DEFAILLANTE': ['HEATER_POMPE_HS', 'NIVEAU_BAS_VANNE_PANNE', 'BULLES_AIR'],
    'BULLES_AIR': ['BULLES_AIR', 'NIVEAU_BAS_VANNE_PANNE'],
    'NIVEAU_BAS': ['NIVEAU_BAS_VANNE_PANNE', 'FUITE_CIRCUIT'],
    'CALCAIRE': ['CALCAIRE_TUYAUX'],
    'NORMAL': [],
    'ERREUR': ['HEATER_POMPE_HS', 'HEATER_RESISTANCE_HS']
}

# -- HEATER TEMPERATURE SENSOR ------------------------------------------------
HEATER_TEMP_SENSOR = None

# -- MODEL HEALTH EVALUATION --------------------------------------------------
EVAL_INTERVAL_CYCLES = 600
EVAL_WINDOW_MINUTES  = 30
EVAL_PERSISTENCE     = 3
IF_ANOMALY_RATE_MAX  = 0.15
RF_F1_WEIGHTED_MIN   = 0.75
RF_CONFIDENCE_THRESHOLD = 0.5
\end{lstlisting}

\endinput
